
\medskip
\phantomsection\label{lem:sga3-VIB-5.6.3}%
\noindent\textbf{Lemma 5.6.3.---}
\emph{Let \(X\) be an \(S\)-scheme. The following conditions are
equivalent:}
\begin{enumerate}
\item[\textup{\SGAInlineRefTarget{lem:sga3-VIB-5.6.3:item:i}(i)}]\phantomsection\label{cond:sga3-VIB-5.6.3-i}
  \emph{\(X\) is separated over \(S\).}

\item[\textup{\SGAInlineRefTarget{lem:sga3-VIB-5.6.3:item:ii}(ii)}]\phantomsection\label{cond:sga3-VIB-5.6.3-ii}
  \emph{For every \(S\)-scheme \(T\), every section
  \(h:T\to X_T\) is a closed immersion.}

\item[\textup{\SGAInlineRefTarget{lem:sga3-VIB-5.6.3:item:iii}(iii)}]\phantomsection\label{cond:sga3-VIB-5.6.3-iii}
  \emph{For every reduced \(S\)-scheme \(T\), two morphisms
  \(f_1,f_2:T\to X\) that coincide on a dense open subset \(U\) of
  \(T\) coincide.}

\item[\textup{\SGAInlineRefTarget{lem:sga3-VIB-5.6.3:item:iv}(iv)}]\phantomsection\label{cond:sga3-VIB-5.6.3-iv}
  \emph{For every \(S\)-scheme \(T\), every point \(t\in T\) and every
  pair of rational points \(x_1,x_2\) of
  \(X_t=X\times_S\operatorname{Spec}\kappa(t)\), there exist a morphism
  \(g:T''\to T'\to T\), with \(T'\to T\) open and \(T''\to T'\) closed and
  dominant, such that \(g^{-1}(t)\ne\varnothing\), and finally an open
  subscheme \(V\) of \(X_{T''}\), separated over \(T''\), which contains
  \(g^{-1}(x_1)\) and \(g^{-1}(x_2)\).}\footnote{Translator's note: the Gille--Polo re-edition (2011),
  Expos\'e~VI\textsubscript{B}, states (iv) for \(T=S\) only, with arbitrary points \(x_1,x_2\) of a fibre
  \(X_s\) and \(V\) containing their inverse image; its proof of (iv)\(\Rightarrow\)(iii) by base change to \(T\)
  then needs \(S'\to S\) universally open and \(S''\to S'\) universally closed and surjective, as they are in the
  application to 5.3 (\'etale, and finite surjective). It adds the condition: (iv\('\)) \emph{every two points of a
  fibre \(X_s\) lie in an open subscheme of \(X\) separated over \(S\)}, which is equivalent to the others, since
  (i) implies (iv\('\)) (take \(X\)), and (iv\('\)) implies (iv) (take \(T''=T'=T\) and the inverse image of such an
  open subscheme).}
\end{enumerate}

\textup{(i)}\(\Rightarrow\)\textup{(iv)}: take \(T''=T\) and \(V=X_T\).

\textup{(iv)}\(\Rightarrow\)\textup{(iii)}: we may suppose \(T=S\). Since \(T\) is reduced, in order to see that
\(f_1=f_2\) it is then enough to see that \(f_1=f_2\) set-theoretically. So let \(t\in T\), and let us show that
\(f_1(t)=f_2(t)\). Let \(g:T''\to T'\to T\) be as in \textup{(iv)}. Since \(U\) is dense in \(T\) and \(T'\to T\)
is open, the inverse image \(U'\) of \(U\) in \(T'\) is dense in \(T'\). Let \(U''\) be the inverse image of
\(U'\) in \(T''\) and \(\overline{U''}\) its closure in \(T''\). Since \(T''\to T'\) is closed and dominant,
\(\overline{U''}\to T'\) is surjective, so \(\overline{U''}\cap g^{-1}(t)\) is not empty. It is then clear that
we may replace \(U\) and \(T\) by the reduced subschemes of \(T''\) having \(U''\) and \(\overline{U''}\) as
underlying spaces, and replace \(f_1\) and \(f_2\) by their inverse images under the base change
\(\overline{U''}\to T\). In short, we may suppose that there exists an open subset \(V\) of \(X_T\), separated
over \(T\), which contains \(f_1(t)\) and \(f_2(t)\). Let \(W=f_1^{-1}(V)\cap f_2^{-1}(V)\); it is an open
subset of \(T\) containing \(t\), and \(U\cap W\) is dense in \(W\). We may therefore replace \(T\) by \(W\). But
then it is clear that \(f_1=f_2\),\footnote{Translator's note: this is the implication
(i)\(\Rightarrow\)(iii) for \(V\), which the Gille--Polo re-edition (2011), Expos\'e~VI\textsubscript{B}, proves;
it also shows why it is enough to prove \(f_1=f_2\) set-theoretically. In our words: if \(X\) is separated over
\(S\) and \(f_1,f_2\colon T\to X\) are \(S\)-morphisms coinciding on a dense open subset of a reduced scheme
\(T\), then \((f_1,f_2)\colon T\to X\times_SX\) maps a dense subset, hence all of \(T\), into the closed diagonal,
and factors through it, \(T\) being reduced; so \(f_1=f_2\). Reducedness is needed: for
\(T=\operatorname{Spec}k[\varepsilon,x]/(\varepsilon^2,\varepsilon x)\) and \(X=\operatorname{Spec}k[x]\), \(k\) a
field, the morphisms \(T\to X\) defined by \(x\mapsto x+\lambda\varepsilon\) \((\lambda\in k)\) coincide on the
dense open subset \(D(x)\), where \(\varepsilon=0\), but are distinct.} hence \(f_1(t)=f_2(t)\).

\textup{(iii)}\(\Rightarrow\)\textup{(ii)}. We may suppose \(T\) reduced. Let \(E\) be the subscheme \(h(T)\) of
\(X_T\), and let \(\overline E\) be the reduced closed subscheme of \(X_T\) having the closure of \(E\) as
underlying space, so that \(E\) is a dense open subscheme of \(\overline E\). After the base change
\(\overline E\to S\), we have two sections of \(X_{\overline E}\) which coincide on \(E\); hence they coincide,
and consequently \(E=\overline E\).

\textup{(ii)}\(\Rightarrow\)\textup{(i)}. Take \(T=X\) and for \(h\) the diagonal section.

The theorem then follows trivially from Lemma \hyperref[lem:sga3-VIB-5.6.3]{5.6.3},
\textup{(iv)}\(\Rightarrow\)\textup{(i)}, and Corollary \hyperref[cor:sga3-VIB-5.6.2]{5.6.2}.
